% Source: https://artofproblemsolving.com/community/c3h1468176_interesting_geometry
% Stable identity: aops-topic-1468176-q01
Let $I$ be the incentre and $O$ the circumcentre, so $O$ is the reflection of $I$ in $AB$.

\textbf{The triangle is isosceles.} $O$ lies on the perpendicular bisector $\ell$ of $AB$. Reflection in line $AB$ maps $\ell$ to itself, so $I$ lies on $\ell$ as well. Then $IA=IB$, so $\angle IAB=\angle IBA$, that is $\dfrac{\angle A}{2}=\dfrac{\angle B}{2}$. Let $\angle A=\angle B=\alpha$.

\textbf{The angle at $C$ is obtuse.} $I$ is inside the triangle, on the same side of $AB$ as $C$, so $O$ is on the other side of $AB$. The circumcentre lies beyond side $AB$ exactly when $\angle C$ is obtuse, so $\angle C=180^\circ-2\alpha>90^\circ$ and $0<\alpha<45^\circ$.

\textbf{Distances to $AB$.} Put $A=(-1,0)$, $B=(1,0)$, $C=(0,\tan\alpha)$. The incentre is on the bisector from $A$, above the midpoint $(0,0)$ of $AB$, so $I=\left(0,\tan\dfrac{\alpha}{2}\right)$. The circumcentre is $O=(0,-d)$ with $OA=OC$:
\[
1+d^2=(\tan\alpha+d)^2\quad\Longrightarrow\quad d=\dfrac{1-\tan^2\alpha}{2\tan\alpha}=\dfrac{1}{\tan2\alpha}=\tan(90^\circ-2\alpha).
\]
The reflection condition is $d=\tan\dfrac{\alpha}{2}$, that is $\tan(90^\circ-2\alpha)=\tan\dfrac{\alpha}{2}$. Both angles lie between $0^\circ$ and $90^\circ$, so $90^\circ-2\alpha=\dfrac{\alpha}{2}$ and $\alpha=36^\circ$.

The angles are $\boxed{\angle A=\angle B=36^\circ,\ \angle C=108^\circ}$.

\begin{asy}
usepackage("amsmath");
size(5.5cm);
defaultpen(fontsize(10pt));
pair A=(-1,0),B=(1,0),C=(0,Tan(36));
real r=Tan(18); pair I=(0,r),O=(0,-r);
real R=length(A-O);
draw(circle(O,R),gray);
draw(circle(I,r),blue);
draw(A--B--C--cycle,linewidth(0.9));
draw(I--O,red+dashed);
dot(A);dot(B);dot(C);dot(I,blue);dot(O,red);
label("$A$",A,W); label("$B$",B,E); label("$C$",C,N);
label("$I$",I,E,blue); label("$O$",O,E,red);
draw(arc(A,0.28,0,36)); draw(arc(B,0.28,144,180));
label("$36^\circ$",A+0.5*dir(13),fontsize(8pt)); label("$36^\circ$",B+0.5*dir(167),fontsize(8pt));
label("\mbox{}",(min(currentpicture,user=true)+max(currentpicture,user=true))/2,black); // keeps the text after the figure black when LaTeX typesets the labels inline
\end{asy}
